\documentclass{homework}
% \usepackage{graphicx} % Required for inserting images
\usepackage{amsmath}
\usepackage{gensymb}
\usepackage{subfig}
\usepackage{float}

\usepackage[framed,numbered,autolinebreaks,useliterate]{mcode}

\class{ECEN 5244}
\title{Homework 2}
\author{Sean Jennings}
\date{September 9, 2024}

\begin{document}
\maketitle

\section{Problem 1}
\begin{letterenum}
\item 
For each ray $i\in\{1, \dots, 5\}$, we are given $d_i$, the travel time of the ray through the material. 
Let $v_j$ be the velocity of the ray as it travels through the $j$th zone where $j\in\{1,\dots,6\}$. From
the given illustration, we write the distance $G_{ij}$ which ray $i$ travels through zone $j$ as
entries of the matrix $G\in\R^{5\times 6}$:
\begin{equation*}
    G = \mat{
        1 & 1 & 1 & 1 & 1 & 1 \\
        \sqrt{2} & 0 & 0 & 0 & 0 & 0 \\
        0 & 0 & 1 & 0 & 0 & 0 \\
        0 & 0 & 0 & \sqrt{2} & 0 & 0 \\
        0 & 0 & 0 & 0 & 0 & 1 \\
    }
\end{equation*}

Assuming constant speed through each zone, the total time spent by ray $i$ in zone $j$ is $G_{ij} / v_j$,
so that the estimated time $\hat{d}_i$ as a function of our velocity model is
\begin{equation*}
    \hat{d}_i = \sum_{j=1}^{6} G_{ij} / v_j
\end{equation*}
Defining $m_j = 1 / v_j$ to be the "slowness" of zone $j$ (the recipricol of velocity), and substituting the slowness into the previous
equation, we see that the travel time is a linear function of the slowness.
\begin{equation*}
    \hat{d}_i = \sum_{j=1}^{6} G_{ij} \cdot m_j
\end{equation*}
We rewrite the previous relationship in vector form with $\hat{d}\in\R^5$ being the vector of estimated travel times for each
ray and $m\in\R^6$ being the modeled slowness in each zone. The relationship is now in our canonical form:
\begin{equation}
    \hat{d} = Gm
    \label{eq:canonical-linear}
\end{equation}

\item The system is underdetermined, since we have more modeled parameters than we have measurements. By inspection of the
$G$ matrix, we see that $G$ has rank 5. Since rank$(G)=$rank$(G^T G)$, we see that $G^T G \in\R^{6 \times 6}$ is not full
rank and therefore is not invertible.

\item The least square problem cannot be solved in its explicit form $\hat{m} = (G^T G)^{-1} G^T d$ because $G^T G$ is not invertible.
Since the range of $G$ is $\R^5$ there are infinitely many $\hat{m}$ which satisfy equation (\ref{eq:canonical-linear}). Possible adjustments which
would make the solution unique include:
\begin{enumerate}
    \item Use SVD to find one particular solution,
    \item Collect more measurements,
    \item Impose constraints on how fast the slowness varies between zones,
    \item Combine zones to reduce the dimensionality of the model.
\end{enumerate}

\end{letterenum}

\section{Problem 2}
Let the matrix $G\in\R^{m \times n}$ be given. Define $A\in \R^{n\times n}$ to be $A = G^T G$.

The given expression $A_{ij} m_i m_j$ is a function of the $i$th and $j$th model parameters so when we 
take the partial derivative in terms of $m_k$, the result depends on the indices:
\begin{equation}
    \frac{\delta}{\delta m_k} A_{ij} m_i m_j =  \biggl\{
        \begin{array}{ll}
            0 &\qquad \text{if } i \neq k \text{ and } j\neq k \\
            A_{kj} m_j &\qquad \text{if } i = k \text{ and } j\neq k \\
            A_{ik} m_i &\qquad \text{if } i \neq k \text{ and } j= k \\
            2A_{kk} m_k &\qquad \text{if } i = j = k
        \end{array}
    \biggr.
    \label{eq:conditional-partial}
\end{equation}

Let $f: \R^n \to \R$ be the function $f(m) = m^T A m$.

To relate the previous conditional equality to the derivative of $f$, we write the matrix product $m^T A m$ as a summation over rows and columns of $A$:
\begin{equation}
    f(m) = m^T A m = \sum_{i=1}{n} m_i \biggl(\sum_{j=1}{n} A_{ij} m_j\biggr) = \sum_{i=1}^{n} \sum_{j=1}^n A_{ij} m_i m_j
    \label{eq:summation-form}
\end{equation}

\renewcommand{\partial}[1]{\frac{\delta}{\delta #1}}

We can then find the partial derivative $f$ with respect to $m_k$ by dividing the summation of (\ref{eq:summation-form}) into terms which correspond
to each conditional of (\ref{eq:conditional-partial}):
\begin{align*}
    \frac{\delta}{\delta m_k} f(m) = \frac{\delta}{\delta m_k} (m^T A m) &=
        \sum_{i\neq k} \biggl(\partial{m_k} ( A_{ik} m_i m_k) + \sum_{j\neq k} \partial{m_k} (A_{ij} m_i m_j) \biggr) \\
        &\quad + \partial{m_k} (A_{kk} m_k m_k) + \sum_{j\neq k} \partial{m_k} (A_{kj} m_k m_j) \\
        &= \sum_{i\neq k} A_{ik} m_i + \sum_{j \neq k} A_{kj} m_j + 2 A_{kk} m_k \\
        &= \sum_{i=1}^{n} A_{ik} m_i + \sum_{j=1}^{n} A_{kj} m_j \\
        &= 2 \sum_{j=1}^{n} A_{kj} m_j
\end{align*}
Note the last equality is due to the synmmetry of $A$. This shows that $\partial{m_k} (m^T A m) = (2Am)_k = 2(G^T Gm)_k$.

\newcommand{\inner}[2]{\langle #1, #2 \rangle}

The author finds the index notation rather cumbersome and will instead provide a proof of the normal equation which is 
based upon the derivative of the inner product. Let $\inner{\cdot}{\cdot}$ be the standard inner product in $\R^m$. That is,
$\inner{x}{y} = y^T x$. If $\alpha$ and $\beta$ are two continuously differentiable functions from $R^n$ to $\R^m$, then
the derivative of their inner product is given by:
\begin{equation}
    \nabla \inner{\alpha(x)}{\beta(x)} = \nabla \alpha(x)^T \beta (x) + \nabla \beta(x)^T \alpha(x)
    \label{eq:product-rule}
\end{equation}
This is an extension of the product rule to vector-valued functions.

If we let $\alpha(m) = \beta(m) = Gm$, then the derivative is constant $\nabla \alpha = \nabla \beta = G$ and equation
(\ref{eq:product-rule}) yields the same inequality shown before: $\nabla (m^T G^T G m) = \nabla (\inner{Gm}{Gm}) = 2 G^T Gm$.

Now, given design matrix $G\in\R^{m\times n}$ and measurement vector $d \in\R^m$, we want to find the model $m\in\R^n$ which minimizies
the square distance between our estimates $\hat{d}(m) = Gm$ and the measurements $d$. Let $r: \R^n \to \R$ be the square distance
function defined by:
\begin{equation*}
    r(m) = \norm{\hat{d}(m) - d}^2 = \norm{Gm - d}^2 = \inner{Gm-d}{Gm-d}
\end{equation*}

Since $r$ is defined over all of $\R^n$, $r$ must take its minimum when $\nabla r(m) = 0$.
Now, let $\alpha(m) = \beta(m) = Gm - d$. Then,
\begin{equation*}
    \alpha(m)^T \beta(m) = G^T( Gm - d)  = \nabla \beta(m)^T \alpha(m)
\end{equation*}

and by equation (\ref{eq:product-rule}), the derivative of $r$ is $\nabla r(m) = 2G^T( Gm - d)$. Therefore, $\nabla r(m) = 0$ implies that
\begin{equation*}
    2G^T(Gm-d) = 0
\end{equation*}

Dividing by 2, distributing $G^T$ and adding $G^Td$ to both sides yields the normal equation:
\begin{equation*}
    G^T G m = G^T d
\end{equation*}


\section{Problem 3}
\begin{letterenum}
\item Figure (\ref{fig:fit1}) shows a plot of the data and fit. Figure (\ref{fig:hist1}) shows a plot of the residuals.

\begin{center}
    \includegraphics[width=0.5\linewidth]{fit_data1.jpg}
    \captionof{figure}{Fit for data1.dat}
    \label{fig:fit1}
\end{center}
\begin{center}
    \includegraphics[width=0.5\linewidth]{Images/histogram_data1.jpg}
    \captionof{figure}{Histogram for data1.dat}
    \label{fig:hist1}
\end{center}

\item The $\chi^2$ statistic is 86.0. Since the degrees of freedom for the distribution is $\nu = N - M = 10 - 2 = 8$,
the expectation is ${\bf E}[\chi^2] = \nu = 8$ and the standard deviation is $\sigma_{\chi^2} = \sqrt{2 \nu} = 4$.

\item Given the assumption that the output follows the fit plus Gaussian noise of std $\sigma= 0.2$, the probability of
realizing such a high $\chi^2$ statistic is near 0 ($3.3\times10^-14$). This our model is not acceptable.

\item Figure (\ref{fig:fit2}) shows a plot of the data and fit for the second data file. Figure (\ref{fig:hist2}) shows a plot of residuals.
\begin{center}
    \includegraphics[width=0.5\linewidth]{Images/fit_data2.jpg}
    \captionof{figure}{Fit for data2.dat}
    \label{fig:fit2}
\end{center}
\begin{center}
    \includegraphics[width=0.5\linewidth]{Images/histogram_data2.jpg}
    \captionof{figure}{Histogram for data2.dat}
    \label{fig:hist2}
\end{center}
\item Since the number of data points and the number of parameters are the same as for the first data file, the expectation and
standard deviation of the $\chi^2$ distribution is equivalent to the previous part. The realized value of $\chi^2$ is now
32.3. Again, this is very far outside of the expected range.

\end{letterenum}

\section{Appendix: Code}


\begin{lstlisting}
function y_fit = ls_fit(x, y)
    M = size(x, 1);
    G = [ones(M, 1), x];
    m = (G' * G) \ (G' * y);
    y_fit = G * m;
end

function chi2 = goodness_of_fit(nu, y, y_fit, sigma)
    chi2 = sum((y - y_fit).^2 / sigma^2);
    p_chi2 = 1 - chi2cdf(chi2, nu);
    
    fprintf('Chi2 is: %g\n', chi2);
    fprintf('Probability of result given model: %g\n', p_chi2);
    fprintf('Degrees of freedom is %g so the std is: %g \n', nu, sqrt(2 * nu));
end

function run_fit(x, y, sigma, data_name)
    IMAGE_DIR = '../../tex/hw2/Images/';

    nu = size(x, 1) - 2;
    y_fit = ls_fit(x, y);
    chi2 = goodness_of_fit(nu, y, y_fit, sigma);

    % Data and fit line
    h = figure;
    plot(x, y, 'ro', x, y_fit, 'b');
    legend('data', 'fit');
    title(['Data and Fit Line for ' data_name])
    xlabel('x')
    ylabel('y')
    saveas(h, [IMAGE_DIR 'fit_' data_name '.jpg']);
    uiwait(h)

    % Residual Histogram
    h = figure;
    histogram(y - y_fit, 10);
    title(['Histogram of Residuals for ' data_name ...
           sprintf(' | goodness of fit = %g', chi2)])
    xlabel('error = y - y_fit')
    ylabel('Count')
    saveas(h, [IMAGE_DIR 'histogram_' data_name '.jpg']);
    uiwait(h)
end

load -ascii data1.dat
disp('Fitting data1.dat')
x = data1(:, 1); y = data1(:, 2);
sigma = 0.2;
run_fit(x, y, sigma, 'data1');

load -ascii data2.dat
disp('Fitting data1.dat')
x = data2(:, 1); y = data2(:, 2);
run_fit(x, y, sigma, 'data2');



\end{lstlisting}
    














\end{document}
